\begin{table}[htbp]
\centering
\caption{Standardized Effect Sizes: Grid Isolation and Economic Outcomes}
\label{tab:sde}
\small
\begin{tabular}{lcccccc}
\hline\hline
Outcome & $\hat{\beta}$ & SE & SD($Y$) & SDE & SE(SDE) & Classification \\
\hline
\\[-6pt]
\multicolumn{7}{l}{\textit{Panel A: Pooled (County-Trend Specification)}} \\[3pt]
Log employment & 0.0153 & 0.0137 & 1.8282 & 0.0084 & 0.0075 & Small positive \\
Log avg weekly wage & 0.0225 & 0.0167 & 0.2559 & 0.0879 & 0.0653 & Moderate positive \\
Log establishments & 0.0032 & 0.0055 & 1.5658 & 0.0020 & 0.0035 & Null \\[6pt]
\multicolumn{7}{l}{\textit{Panel B: Heterogeneous (Sample Splits)}} \\[3pt]
Log emp (Q1 2021 only) & 0.0008 & 0.0123 & 1.8282 & 0.0004 & 0.0067 & Null \\
Log emp (rural only) & 0.0148 & 0.0142 & 1.5684 & 0.0094 & 0.0091 & Small positive \\
\hline\hline
\end{tabular}
\begin{tablenotes}
\small
\item \textit{Notes:} \textbf{Country:} United States. \textbf{Research question:} Does electrical grid isolation from the national interconnection impose economic costs when extreme weather causes infrastructure failure, as measured by county-level employment, wages, and business establishments? \textbf{Policy mechanism:} ERCOT operates Texas's isolated electrical grid, disconnected from national interconnections to avoid federal regulation; during Winter Storm Uri (February 2021), this isolation prevented power imports during cascading generator failures, causing 4.5 million customers to lose electricity for up to 5 days. \textbf{Outcome definition:} Log quarterly private-sector employment from BLS QCEW (ownership code 5), measuring average of three monthly employment levels per quarter. \textbf{Treatment:} Binary --- county served by ERCOT (isolated grid) vs.\ SPP/MISO/WECC (nationally connected grids). \textbf{Data:} BLS Quarterly Census of Employment and Wages, 254 Texas counties, Q1 2018--Q4 2023, 6,096 county-quarter observations. \textbf{Method:} Two-way fixed effects DiD with county and year-quarter fixed effects, county-specific linear trends (preferred specification); standard errors clustered at the county level. \textbf{Sample:} All 254 Texas counties with non-suppressed QCEW data; 213 ERCOT, 41 non-ERCOT (26 SPP, 13 MISO, 2 WECC). SDE $= \hat{\beta} / \text{SD}(Y)$ where SD($Y$) is the pre-treatment standard deviation of the log outcome. Classification refers to magnitude, not statistical significance: Large ($|$SDE$| > 0.15$), Moderate ($0.05$--$0.15$), Small ($0.005$--$0.05$), Null ($< 0.005$).
\end{tablenotes}
\end{table}

